\documentclass[11pt]{article}
\usepackage[margin=1in]{geometry}
\usepackage{booktabs}
\usepackage{enumitem}
\usepackage{hyperref}
\usepackage{amsmath}
\usepackage{xcolor}

\title{Revised Project Plan:\\Closed-Loop EEG Cursor Control via Constrained RL\\[0.3em]
\large Offline-to-Online Simulation with Comparative Decoder--Policy Analysis}
\author{BCI Course Project --- Spring 2026}
\date{April 2026}

\begin{document}
\maketitle

%--------------------------------------------------------------
\section{Project Summary}
We build a closed-loop BCI cursor-control pipeline in simulation.
A frozen neural decoder (EEGNet or LSTM) maps 64-channel motor-imagery EEG to 2-D cursor velocity.
A Proximal Policy Optimization (PPO) agent sits on top of the decoder and outputs a \emph{corrective velocity offset} to improve target acquisition under realistic noninvasive latency (500--2000\,ms).
The core comparison is \textbf{Naive PPO vs.\ Constrained PPO} (smoothness + zeroness penalties), evaluated across two decoder architectures.

%--------------------------------------------------------------
\section{Identified Gaps in Original Proposal}

\begin{enumerate}[leftmargin=*]
    \item \textbf{Latency range unrealistic.} Original 50--150\,ms matches intracortical BCIs; noninvasive EEG decoders exhibit 0.5--2\,s lag~\cite{korik2024}.
    \item \textbf{No biomechanical constraints on RL agent.} A simple jerk penalty ($\lambda\|\Delta v\|$) allows unphysical solutions; smoothness and zeroness constraints are needed~\cite{pitt2024}.
    \item \textbf{SAC vs.\ PPO inconsistency.} Prior BCI--RL work uses PPO; we standardize on PPO for comparability.
    \item \textbf{No curriculum learning.} RL training on small targets from scratch is unstable; curriculum over target radius is essential.
    \item \textbf{Missing evaluation metrics.} Added First-Touch Time, Dial-In Time, Distance Ratio, and Shannon--Welford model fit.
    \item \textbf{No non-RL adaptation baseline.} Added Behavioral Cloning (BC) for a supervised-learning reference.
\end{enumerate}

%--------------------------------------------------------------
\section{Architecture}

\begin{verbatim}
Raw EEG (64ch, 512Hz)
  -> Preprocessing (bandpass, notch, CAR, ICA, epoch)
    -> Frozen Decoder (EEGNet / LSTM) -> decoded velocity
                                              |
      cursor position + target position ------+
                                              v
                                    PPO Agent (corrective offset)
                                              |
                                              v
                                    Final cursor velocity
                                              |
                                     Latency buffer (tau ~ U[500,2000] ms)
                                              |
                                              v
                                    Simulation environment (Gym)
\end{verbatim}

%--------------------------------------------------------------
\section{Comparative Dimensions}

\begin{table}[h]
\centering
\caption{Methods compared along each dimension.}
\begin{tabular}{@{}lll@{}}
\toprule
\textbf{Dimension} & \textbf{Methods} & \textbf{Rationale} \\
\midrule
Decoder        & EEGNet vs.\ LSTM RNN       & Shallow CNN vs.\ recurrent temporal model \\
RL Policy      & Naive PPO vs.\ Constrained PPO & Main contribution: effect of constraints \\
Baseline       & LDA decoder (no RL)        & Classical linear baseline \\
Baseline       & Behavioral Cloning (BC)    & Non-RL adaptation reference \\
\bottomrule
\end{tabular}
\end{table}

%--------------------------------------------------------------
\section{Reward Function}

The Constrained PPO reward at time step $t$:
\begin{equation}
r_t = -\underbrace{\|p_t - g_t\|_2}_{\text{distance to target}}
      \;-\; \lambda_{\text{jerk}}\underbrace{\|\Delta v_t\|_2}_{\text{jerk penalty}}
      \;-\; \lambda_{\text{smooth}}\underbrace{D_{\mathrm{KL}}(\pi_t \| \pi_{t-1})}_{\text{smoothness}}
      \;-\; \lambda_{\text{zero}}\underbrace{\|v_t\| \cdot \mathbf{1}[\|p_t - g_t\| < \epsilon]}_{\text{zeroness}}
\end{equation}
where $\lambda_{\text{jerk}}=0.1$, $\lambda_{\text{smooth}}\in\{0.05,0.1\}$, $\lambda_{\text{zero}}\in\{0.05,0.1\}$, and $\epsilon$ is the target acceptance radius.
Naive PPO uses only the first two terms.

%--------------------------------------------------------------
\section{Six-Week Sprint Plan}

\subsection*{Week 1: Data + Preprocessing + Baselines}
\begin{itemize}[leftmargin=*,nosep]
    \item Download Continuous Pursuit EEG dataset~\cite{korik2024} (28 subjects, 64-ch, 512\,Hz).
    \item Preprocessing: bandpass 1--40\,Hz (4th-order Butterworth), 50\,Hz notch, common-average reference, ICA artifact rejection (MNE-Python), epoch into 500\,ms windows with 250\,ms stride, $z$-score per channel per session.
    \item Session-level splits: Sessions 1--2 $\to$ train (70\%), Session 3 $\to$ RL fine-tuning (15\%), Session 4+ $\to$ test (15\%).
    \item Train LDA baseline decoder; record baseline TAR and NMSE.
    \item \textbf{Deliverable:} Clean data pipeline, LDA baseline numbers.
\end{itemize}

\subsection*{Week 2: Offline Decoders (EEGNet + LSTM)}
\begin{itemize}[leftmargin=*,nosep]
    \item Implement EEGNet (depthwise separable CNN) $\to$ 2-D velocity.
    \item Implement LSTM decoder (same input/output interface).
    \item Training: MSE + $\ell_2$ regularization, Adam ($\eta=10^{-3}$, cosine decay), batch size 64, early stopping on validation NMSE.
    \item Evaluate both on held-out sessions: TAR, NMSE, Path Efficiency.
    \item \textbf{Deliverable:} Two trained decoders with offline comparison table.
\end{itemize}

\subsection*{Week 3: Simulation Environment + RL Setup}
\begin{itemize}[leftmargin=*,nosep]
    \item Build Gymnasium-compatible environment:
    \begin{itemize}[nosep]
        \item \textbf{State:} decoder output $\oplus$ cursor position $\oplus$ target position $\oplus$ velocity history (last 5 steps).
        \item \textbf{Action:} corrective velocity offset (continuous, clipped to $[-1,1]$).
        \item \textbf{Latency:} stochastic delay $\tau \sim \mathcal{U}[500, 2000]$\,ms.
    \end{itemize}
    \item Implement curriculum learning: target acceptance radius $120\,\text{mm} \to 60\,\text{mm} \to 30\,\text{mm}$.
    \item Freeze decoder weights throughout RL training.
    \item \textbf{Deliverable:} Working Gym environment with latency injection and curriculum.
\end{itemize}

\subsection*{Week 4: RL Training + BC Baseline}
\begin{itemize}[leftmargin=*,nosep]
    \item \textbf{Naive PPO:} Stable-Baselines3 PPO, reward uses distance + jerk terms only.
    \item \textbf{Constrained PPO:} Add smoothness (KL) and zeroness penalties to reward.
    \item \textbf{Behavioral Cloning:} Supervised policy trained on best offline trajectories.
    \item Hyperparameter grid (12 configs): lr $\in\{3\!\times\!10^{-4}, 1\!\times\!10^{-4}\}$, $\lambda_{\text{smooth}}\in\{0.05, 0.1\}$, $\lambda_{\text{zero}}\in\{0.05, 0.1\}$.
    \item Train on Session 3 data from 5 subjects.
    \item \textbf{Deliverable:} Trained policies --- 4 conditions per decoder (No-RL, BC, NaivePPO, ConstrainedPPO).
\end{itemize}

\subsection*{Week 5: Evaluation + Cross-Subject Generalization}
\begin{itemize}[leftmargin=*,nosep]
    \item Full metric suite on held-out sessions:
    \begin{itemize}[nosep]
        \item \textbf{Primary:} TAR ($>$55\%), NMSE ($<$0.30), Path Efficiency ($>$0.60), Jerk reduction ($>$30\%).
        \item \textbf{Secondary:} First-Touch Time, Dial-In Time, Distance Ratio, Bitrate.
        \item \textbf{Bonus:} Shannon--Welford model fit vs.\ Fitts's Law.
    \end{itemize}
    \item Cross-subject: train on 23 subjects, test on 5 held-out.
    \item Statistical testing: paired Wilcoxon signed-rank, Bonferroni correction ($p<0.05$).
    \item \textbf{Deliverable:} Results tables, trajectory visualizations, statistical analysis.
\end{itemize}

\subsection*{Week 6: Report + Deliverables}
\begin{itemize}[leftmargin=*,nosep]
    \item IEEE 8-page report:
    \begin{itemize}[nosep]
        \item Main comparison table: LDA vs.\ EEGNet vs.\ LSTM $\times$ NoRL vs.\ BC vs.\ NaivePPO vs.\ ConstrainedPPO.
        \item Trajectory plots (best/worst subjects).
        \item Ablation: effect of latency range, curriculum, constraints.
    \end{itemize}
    \item Code freeze, Docker image, reproducibility check (fixed seeds, W\&B logging).
    \item Poster draft (A0) + 5-min video demo.
    \item \textbf{Deliverable:} Final submission package.
\end{itemize}

%--------------------------------------------------------------
\section{Evaluation Summary}

\begin{table}[h]
\centering
\caption{Evaluation metrics and targets.}
\begin{tabular}{@{}llc@{}}
\toprule
\textbf{Metric} & \textbf{Description} & \textbf{Target} \\
\midrule
Target Acquisition Rate (TAR) & \% trials cursor reaches target & $>$55\% \\
Normalized MSE (NMSE) & Cursor--target trajectory error & $<$0.30 \\
Path Efficiency (PE) & Ideal / actual path length & $>$0.60 \\
Jerk Reduction & $\|\ddot{v}\|_2$ mean vs.\ offline & $>$30\% \\
First-Touch Time (FTT) & Time to first target entry & Lower is better \\
Dial-In Time (DIT) & Time to stabilize on target & Lower is better \\
Distance Ratio & Actual / straight-line path & Closer to 1.0 \\
Inference Latency & Decoder + policy wall-clock & $<$50\,ms \\
\bottomrule
\end{tabular}
\end{table}

%--------------------------------------------------------------
\section{Resources}

\begin{itemize}[leftmargin=*,nosep]
    \item \textbf{Compute:} 1$\times$ NVIDIA A100 40\,GB (university HPC); est.\ 48 GPU-hours.
    \item \textbf{Libraries:} Python 3.11, MNE-Python 1.7, PyTorch 2.3, Stable-Baselines3 2.3, Gymnasium, NumPy/SciPy.
    \item \textbf{Data:} Continuous Pursuit EEG BCI Dataset~\cite{korik2024} (CC-BY 4.0).
    \item \textbf{Reproducibility:} Fixed seeds (NumPy, PyTorch, SB3), DVC for data versioning, Docker image, W\&B tracking.
\end{itemize}

%--------------------------------------------------------------
\begin{thebibliography}{9}
\bibitem{korik2024} A.~Korik et al., ``A continuous pursuit dataset for online deep learning-based EEG brain-computer interface,'' \textit{Scientific Data}, vol.~11, 2024.
\bibitem{pitt2024} B.~Pitt and R.~Kramer, ``A reinforcement learning based software simulator for motor brain-computer interfaces,'' \textit{PMC}, 2024.
\bibitem{shin2022} N.~Shin et al., ``Closed-loop motor imagery EEG simulation for brain-computer interfaces,'' \textit{Frontiers in Human Neuroscience}, vol.~16, 2022.
\bibitem{lawhern2018} V.~J.~Lawhern et al., ``EEGNet: A compact convolutional neural network for EEG-based brain-computer interfaces,'' \textit{J.\ Neural Eng.}, vol.~15, no.~5, 2018.
\end{thebibliography}

\end{document}
